\subsection{形式幂级数环的因子分解性}

命题 8. 设 C 是一个环，它或者是域，或者是离散赋值环。则形式幂级数整环 C{[}{[}X₁, . . ., Xₙ{]}{]} 是因子分解整环。

令 p 为 C 的极大理想，\(\pi\) 为 p 的一个生成元（若 C 是域，则 \(\pi = 0\)）。赋予 C p-进拓扑，该拓扑是 Hausdorff 的。由于 C 是 Noether 局部环，\(B = C[[X_{1}, \ldots, X_{n}]]\) 是 Noether 局部环，且其完备化为 \(\hat{C}[[X_{1}, \ldots, X_{n}]]\)（第三章，§2，第 6 节，命题 6）。由命题 4 的推论（第 6 节），只需证明 \(\hat{C}[[X_{1}, \ldots, X_{n}]]\) 是因子分解整环。现在，如果 C 是域，则 \(\hat{C} = C\)；如果 C 是离散赋值环，\(\hat{C}\) 也具有同样性质（第六章，§5，第 3 节，命题 5）。因此在证明的其余部分，假设 C 是完备的。

从 n = 0 的平凡情形开始作归纳，假设已经证明 \(A = C[[X_{1}, \ldots, X_{n-1}]]\) 是因子分解整环。我们把 B 识别为 \(A[[X_{n}]]\)，并以 m 表示 A 的极大理想（由 \(\pi\)、\(X_{1}, \ldots, X_{n-1}\) 生成）。我们将证明，B 的每个非零元 g 都能以本质唯一的方式写成极端元的乘积。

令 K 为域 C/Cπ；由于 B/Bπ 识别为 K{[}{[}X₁, . . . , Xₙ{]}{]}，理想 Bπ 是素理想，而 π 是极端元。若 π ≠ 0，则 \(B_{B\pi}\) 是一个赋范离散赋值 w 的环（第六章，§ 3，第 6 节，命题 9）；因此 B 的每个非零元 g 都可写成 g = πw(g)f，其中 f ∈ B 且 f 不是 π 的倍元。因此只需证明 f 是极端元的本质唯一乘积。现在 f 在 K{[}{[}X₁, . . . , Xₙ{]}{]} 中的典范像非零；所以引理 3（第 7 节）表明，存在 B 的一个自同构，将 f 映为一个元 f'，使得 f'(0, . . . , 0, Xₙ) 的系数不全属于 Cπ；这意味着，把级数 f' 看作关于 Xₙ 的形式幂级数时，其系数不全属于 m。只需对 f' 证明我们的断言。

下文中，B 的所有元素都看作系数在 A 中的 \(X_{n}\) 的形式幂级数。由第 8 节命题 6（之所以适用，是因为 C，从而 A，是分离且完备的，并且 \(f'\) 的约化幂级数 \(\neq0\)），\(f'\) 在 B 中与唯一的特异多项式 F 相伴。由第 8 节命题 7，任何整除 \(f'\) 的级数（或者等价地，整除 F 的级数）都与一个整除 F 的特异多项式相伴，而 \(f'\) 的每个分解，在可逆因子意义下，都形如 \(f' = uF_{1}\ldots F_{q}\)，其中 u 可逆，且 \(F_{t}\) 是 B 中满足 \(F = F_{1}\ldots F_{q}\) 的极端特异多项式。由第 8 节命题 7 的推论，\(F_{t}\) 在 \(A[X_{n}]\) 中也都是极端元。现在，由于 A 按归纳假设是因子分解整环，\(A[X_{n}]\) 也是如此（定理 2，第 5 节）；因此，由于它们是首一的，\(F_{i}\) 由 F 唯一确定（相差一个排列）。这说明分解 \(f' = uF_{1} \ldots F_{q}\) 的唯一性；其存在性来自 B 是 Noether 环这一事实，证明完毕。

\textbf{注}\par

\begin{enumerate}
\def\labelenumi{(\arabic{enumi})}
\item
  存在因子分解整环 A，使得环 A{[}{[}X{]}{]} 不是因子分解整环（习题 8）。然而，若 A 是主理想整环，则 A{[}{[}X₁, \ldots, Xₙ{]}{]} 是因子分解整环（习题 9）。
\item
  * 我们稍后将用同调方法看到，每个正则局部环都是因子分解整环（参见 § 4，第 7 节，命题 16 的推论 3）。这将给出命题 8 的另一种、在概念上更简单的证明。 *
\end{enumerate}

